\chapter{Library functional}\label{app:functional}

The \snippet{functional} library from tabularray is very powerful. It can modify cell contents, their style, etc., before the table is displayed. However, it relies on the \snippet{functional} package, which itself is a modification of the Latex3 syntax, and thus not the most beginner friendly. Nevertheless, after a few dark magic incantations, the results can be stunning. 

Here is an example, that automates the tikz drawing of the table from \Cref{sec:libraries}. Do not panic if you do not understand the code, just remember that the library exists \emoji{slightly-smiling-face}. 

Our preamble gets some weird voodoo:

\begin{Code}
\documentclass{article}  % Or any other documentclass
\usepackage{tabularray}
\UseTblrLibrary{functional, tikz}
\usepackage{xcolor}

% Inspired from https://tex.stackexchange.com/a/750321/430417
\IgnoreSpacesOn
% Variable must be def before the function
\tlNew \gTikzPercentagePathsTl

% Computes a column's maximum.
% Useful if numbers are not ordered.
\prgNewFunction \colMaxValue {m} {  % takes a mandatory argument, #1, the column where we want to compute the max
	\fpSet \lTmpaFp {\fpEval {\cellGetText{2}{#1}}}  % saving temp max in \lTmpaFp 
	\intStepOneInline {2} {\arabic{rowcount}} {  % ##1 starts from 2 until the number of rows
		\fpSet \lTmpaFp {\fpMathMax{\lTmpaFp}{\fpEval {\cellGetText{##1}{#1}}}}  % update temp max
	}
	\prgReturn {\fpUse\lTmpaFp}
}

\prgNewFunction \barPlotPaths {m} {  % #1 is the index of the column we want the bars
	\tlClear \gTikzPercentagePathsTl  % set the var in which we will store the drawing instructions
	\fpSet \lTmpaFp {\colMaxValue{#1}}  % Set #1 column max value in \lTmpaFp
	\intStepOneInline {2} {\arabic{rowcount}} {  % ##1 starts from 2 until max row
		\fpSet \lTmpbFp {\fpEval{\cellGetText{##1}{#1}}}  % save cell (##1, #1) value into \lTmpbFp 
		\fpSet \lTmpcFp {\fpMathDiv {\lTmpbFp} {\lTmpaFp}}  % save (cell value) / (max value) in \lTmpcFp
		\tlPutRight \gTikzPercentagePathsTl {  % Update drawing instruction with a rectangle of size \lTmpcFp
			\evalWhole{
				(##1-#1.south~west) rectangle
				($(##1-#1.north~west)!{\fpUse \lTmpcFp}!(##1-#1.north~east)$)
			}
		}
	}
}
\IgnoreSpacesOff
\end{Code}

So that our table does not require any copy pasting anymore!

\begin{Code}[show-output]
\begin{tblrtikzbelow}
	% draw everything from \gTikzPercentagePathsTl
	\highLight{\draw[white, line width=0.5pt, fill=cyan!20] \gTikzPercentagePathsTl;}
\end{tblrtikzbelow}

% Original table and design from https://www.storytellingwithdata.com/blog/2012/02/grables-and-taphs
\begin{tblr}{
	colspec={lcccQ[2.5cm]},
	row{1} = {c, m, font=\bfseries, gray9},
	\highLight{process={\barPlotPaths{5}},}  % We call the function barPlotPath defined in the preamble to fill \gTikzPercentagePathsTl with bars to plot in the 5th column
}
	Borough & {Trust \\ rank} & {Index \\ rank } & {Number \\ of grants} & Amount approved (£) \\
	% no more data copy-pasting!
	Tower Hamlets          & 1  & 3  & 269 & 9692642 \\
	Hackney                & 2  & 2  & 225 & 7809608 \\
	Southwark              & 3  & 12 & 232 & 7266118 \\
	Camden                 & 4  & 14 & 136 & 6140419 \\
	Islington              & 5  & 4  & 156 & 5424137 \\
	Lambeth                & 6  & 8  & 156 & 5257941 \\
	Newham                 & 7  & 2  & 154 & 5217075 \\
	Hammersmith and Fulham & 8  & 13 & 109 & 4085708 \\
\end{tblr}
\end{Code}
